\section{Primary-source locator ledger}
\label{app:source-ledger}

The proof uses external sources only for general statements.  Every
instance-specific polynomial, factorization, group count, and matrix rank is
supplied by the exact C56 replay.

\begin{longtable}{@{}>{\raggedright\arraybackslash}p{0.18\textwidth}>{\raggedright\arraybackslash}p{0.29\textwidth}>{\raggedright\arraybackslash}p{0.43\textwidth}@{}}
\toprule
source & exact locator & use and boundary\\
\midrule
\endfirsthead
\toprule
source & exact locator & use and boundary\\
\midrule
\endhead
Kass--Wickelgren
\cite{KassWickelgren2021}
& Theorem~2; \S5, Definition~41; Corollary~53, p.~701;
Corollary~54, p.~702
& Definition~41 identifies the zero scheme of the
\(\operatorname{Sym}^3(\cS^\vee)\) section with the line scheme.
Theorem~2 has rank \(27\).  Corollary~53 gives simple zeros on a smooth
cubic, and Corollary~54 gives separable residue fields.  Corollary~53 alone
is not cited for the number \(27\).\\
\addlinespace
Elsenhans--Jahnel
\cite{ElsenhansJahnel2009}
& Fact~3; Remarks~4--5; Lemma~8; Algorithm~10;
Remarks~11--13; author-PDF pp.~3--7
& Fact~3 places the common line-field group in \(\WE\).  Lemma~8 reduces a
transitive subgroup containing order five to \(\Ucox\) or \(\WE\).
Remarks~11--13 certify the role of type \((2,5,5,5,10)\).
Remark~5 says the \(27\)-line action lies in \(A_{27}\); ordinary
\(S_{27}\) sign does not define \(\Ucox\).\\
\addlinespace
Viray \cite{Viray2023}
& \S2.1, p.~10, equation~(24) and the immediately following low-degree
exact sequence
& Supplies the Picard-to-Brauer segment of Hochschild--Serre.  Since the
field Brauer group is torsion, it yields rank equality after tensoring with
\(\Q\), not integral surjectivity and not a Brauer--Manin conclusion.\\
\addlinespace
Das \cite{Das2021}
& Introduction and the incidence cover
\(\widetilde{\mathcal M}\to\mathcal M\)
& Context for the \(27\)-sheeted moduli cover that marks one line.  It does
not determine the arithmetic fibre studied here.\\
\addlinespace
Pichon-Pharabod--Telen
\cite{PichonPharabodTelen2025}
& Abstract and \S1 around Theorem~2 and Table~1
& Recent context showing that symmetric cubic-surface families often have
proper monodromy subgroups.  It is not an exact line-field calculation for
\cref{eq:cubic}.\\
\addlinespace
McKean \cite{McKean2021}
& Theorems~1.1--1.2, \S1.2, and Appendix~A
& Context for rational-line counts and subgroup actions on the Schläfli
graph.  The no-line conclusion here instead follows from
\(\FanoY=\Spec(E)\).\\
\bottomrule
\end{longtable}

The Elsenhans--Jahnel chapter belongs to Progress in Mathematics
volume~269; the publisher record carries copyright year 2009 and electronic
publication in 2010.  The bibliography uses the volume year 2009.  The
remaining published metadata were checked against the DOI records, and the
preprint metadata against the official arXiv records, on 2026-08-15.

The recent-neighbor search used the query families listed in
\cref{sec:context-scope}.  No temporary download path or digest is a
bibliographic or theorem authority.
